% from ega0-19-fr.tex
\oldpage[0\textsubscript{IV-1}]{77}

$A_{\mathfrak p}$-module et en notant que (puisque $N$ est libre de type fini)
toute forme linéaire sur le $A_{\mathfrak p}$-module $N_{\mathfrak p}$ s'écrit
$u=v/f$, où $f\in A-\mathfrak p$, et où $v$ est une forme linéaire sur le
$A$-module $N$. Il est clair que c) entraîne b), et il reste donc à voir que a)
entraîne d). Or, comme $N$ est de présentation finie,
$(\operatorname{Hom}_A(N,M))_{\mathfrak p}$ s'identifie canoniquement à
$\operatorname{Hom}_{A_{\mathfrak p}}(N_{\mathfrak p},M_{\mathfrak p})$
(Bourbaki, \emph{Alg. comm.}, chap.~II, §~2, n\textsuperscript{o}~7,
prop.~19). Si $w_{\mathfrak p}$ est inverse à gauche de $u_{\mathfrak p}$, il
existe donc un homomorphisme $w:N\to M$ et un élément $f\in A-\mathfrak p$
tels que $w=w_{\mathfrak p}\otimes(1/f)1_{A_{\mathfrak p}}$. La relation
$w_{\mathfrak p}\circ u_{\mathfrak p}=1_{M_{\mathfrak p}}$ s'écrit donc aussi
$(w\circ u)\otimes1_{A_{\mathfrak p}}=f.1_{M_{\mathfrak p}}$. Mais comme $M$
est un $A$-module de type fini, il existe $g\in A-\mathfrak p$ tel que
l'endomorphisme $g((w\circ u)-f.1_M)$ s'annule en tous les générateurs de $M$,
donc soit nul. Si on pose $h=fg$, et
$u_h=u\otimes1_{A_h}$, $w_h=w\otimes1_{A_h}$, on a donc
$w_h(u_h(z))=(f/1)z$ pour tout $z\in M_h$~; mais comme $h/1$ est inversible
dans $A_h$, il en est de même de $f/1=f'$, et $f'^{-1}w_h$ est par suite
inverse à gauche de $u_h$, ce qui achève la démonstration.

\begin{proposition}[19.1.13]
\label{0.19.1.13-fr}
Supposons vérifiées les hypothèses de (19.1.9) et supposons en outre que
$(\mathfrak J_\lambda)$ soit une suite décroissante $(\mathfrak J_n)$ et que
$M$ soit séparé et complet. Alors les conditions a) et b) de (19.1.9) sont
aussi équivalentes à~:
\begin{enumerate}
  \item[c)] $u$ est inversible à gauche.
\end{enumerate}
\end{proposition}

On sait déjà que c) entraîne a) (19.1.6). Inversement, si a) est vérifiée, on
sait (avec les notations de (19.1.8)) que $M_n$ est facteur direct du
$(A/\mathfrak J_n)$-module projectif $N_n$, donc $P_n$ est aussi isomorphe à
un facteur direct de $N_n$, et par suite est projectif~; il suffit donc
d'appliquer (19.1.8).

Notons enfin la proposition suivante~:

\begin{proposition}[19.1.14]
\label{0.19.1.14-fr}
Soient $A$ un anneau, $M$ un $A$-module de type fini, $N$ un $A$-module
projectif, $u:M\to N$ un $A$-homomorphisme.
\begin{enumerate}
  \item[(i)] Pour que $u$ soit inversible à gauche, il faut et il suffit que,
  pour tout idéal maximal $\mathfrak m$ de $A$,
  $u_{\mathfrak m}:M_{\mathfrak m}\to N_{\mathfrak m}$ soit inversible à
  gauche.
  \item[(ii)] Soit $A'$ une $A$-algèbre qui soit un $A$-module fidèlement
  plat. Pour que $u$ soit inversible à gauche, il faut et il suffit que
  $u\otimes1:M\otimes_AA'\to N\otimes_AA'$ soit inversible à gauche.
\end{enumerate}
\end{proposition}

Comme dans (19.1.12), on peut se borner au cas où $N$ est libre de type fini~;
dire que $u$ est inversible à gauche signifie alors que $u$ est injectif et
que le module quotient $P=N/u(M)$ est \emph{projectif}, car $u(M)$ sera alors
facteur direct de $N$. Notons en outre que puisque $M$ est de type fini, $P$
est de \emph{présentation finie}. Cela étant~:

(i) La condition est évidemment nécessaire. Inversement, si elle est
satisfaite, on sait que $u$ est injectif (Bourbaki, \emph{Alg. comm.},
chap.~II, §~3, n\textsuperscript{o}~3, th.~1) et comme
$P_{\mathfrak m}=N_{\mathfrak m}/u_{\mathfrak m}(M_{\mathfrak m})$ est
projectif pour tout $\mathfrak m$, on sait que cela entraîne que $P$ est
projectif (\emph{loc. cit.}, §~5, n\textsuperscript{o}~2, th.~1).

(ii) Ici encore, la condition est trivialement nécessaire. Inversement, si
elle est remplie, on sait que $u$ est injectif ($\mathbf{0}_{\mathrm I}$,
6.4.1) et comme $P\otimes_AA'=\operatorname{Coker}(u\otimes1)$ est projectif,
donc plat, on en déduit que $P$ est un $A$-module plat
($\mathbf{0}_{\mathrm I}$, 6.6.3), donc projectif puisqu'il est de présentation
finie (Bourbaki, \emph{Alg. comm.}, chap.~II, §~5,
n\textsuperscript{o}~2, cor.~2 du th.~1).

